\textbf{Neural and kernel potentials.} Behler and Parrinello introduced neural-network potentials in which the energy is a sum of atomic contributions computed from atom-centred symmetry functions of the local environment~\cite{behler2007generalized}; the radial and angular function families and their cutoff are detailed in~\cite{behler2011atom}, and we use that construction. The summed-atomic-energy form is an instance of the sum-decomposition that characterises permutation-invariant functions on sets~\cite{zaheer2017deep}. Bart\'ok et al.\ review representations of atomic neighbourhoods and propose the SOAP similarity for kernel regression~\cite{bartk2012representing}; we use a plain RBF kernel on symmetry-function vectors instead. Barthelm\'e et al.\ show that Gaussian-kernel regression tends to polynomial regression in the flat limit of very wide kernels, and that the best predictions may lie at length-scales that are numerically hard to reach~\cite{barthelm2022gaussian}; our symmetry-function KRR turned out to sit in this regime.

\textbf{Ordering-free global representations.} Rupp et al.\ encode molecules with a Coulomb matrix and remove its dependence on atom ordering by using the sorted eigenvalue spectrum as the descriptor for kernel ridge regression~\cite{rupp2011fast}. Our sorted-distance descriptor follows the same idea of a canonical ordering, applied directly to the list of inverse pair distances of a single-element cluster.

\textbf{Comparisons and reviews.} Zuo et al.\ compare descriptors and regressors for interatomic potentials in terms of accuracy and cost~\cite{zuo2019performance}, and Parsaeifard et al.\ assess the structural resolution of fingerprints used in machine learning of atomistic systems~\cite{parsaeifard2020assessment}; a broad review of machine-learned force fields is~\cite{unke2020machine}. SchNet~\cite{schtt2017schnet} and NequIP~\cite{batzner2021e} learn invariant or equivariant features end to end and are out of scope here. The minima of LJ clusters are tabulated in~\cite{wales1997global}. Relative to these works we offer only a small reproducible comparison, not a new representation.
